Previous work has applied ethnographic \citep{udupa_prompt_2026} and critical domain analysis \citep{kopf_content_2019, kopf_corporate_2024} close-reading frameworks to a variety of corporate-backed sociotechnical documents such terms of service guidelines and generative AI system prompts, but to our knowledge such frameworks have not yet been turned towards community-generated prompts or artifacts, nor towards the communities themselves. In particular, existing critical methodologies for close reading of sociotechnical documents do not center  community-created and community-spread prompts as artifacts of \textit{knowledge and practice transfer} within harmful contexts. Our work seeks to fill this gap by applying qualitative content analysis using frameworks derived from interpersonal violence scholarship \citep{stark_coercive_2023} and structural examination of pedophile manuals \cite{blalock_content_2022}.